% Source: https://artofproblemsolving.com/community/c3h1676792_35_let_n_be_a_nonnegative_integer_prove_that_sum_i0k_1ibinomn
% Stable identity: aops-topic-1676792-q01
Pascal's rule says $\dbinom ni=\dbinom{n-1}{i}+\dbinom{n-1}{i-1}$ for $i\ge1$ (valid for every $n\ge1$, also when $i>n$ with the convention above), and $\dbinom n0=\dbinom{n-1}{0}=1$. Write $c_i=(-1)^i\dbinom{n-1}{i}$. Then for $i\ge1$
\[
(-1)^i\dbinom ni=(-1)^i\dbinom{n-1}{i}+(-1)^i\dbinom{n-1}{i-1}=c_i-c_{i-1},
\]
and for $i=0$ the term is $1=c_0$. The sum telescopes:
\[
\sum_{i=0}^{k}(-1)^i\dbinom ni=c_0+(c_1-c_0)+(c_2-c_1)+\cdots+(c_k-c_{k-1})=c_k=(-1)^k\dbinom{n-1}{k}. \qquad\blacksquare
\]
For example, $n=5$, $k=2$: $1-5+10=6=\dbinom42$. (For $k\ge n$ both sides are $0$, since a full alternating row of Pascal's triangle sums to $(1-1)^n=0$.)

The hypothesis $n\ge1$ is needed: for $n=0$ and $k=1$ the left side is $\dbinom00-\dbinom01=1$, while the right side $-\dbinom{-1}{1}$ is either undefined or (with the convention above) $0$.
